%%
%% This is a placeholder manuscript shell for EMNLP 2027.
%% Refresh this manuscript before submission by syncing with the research repository upon activation.
%%
\documentclass[11pt]{article}
\usepackage[review]{acl}
\usepackage{times}
\usepackage{latexsym}
\usepackage{microtype}
\usepackage{booktabs}
\usepackage{graphicx}
\usepackage{amsmath}

\title{Empirical Evaluation of Cross-Lingual Prompt Injection: Calibration Disparity and Tokenisation Drift}

\author{Anonymous EMNLP Submission}

\begin{document}
\maketitle

\begin{abstract}
Empirical evaluation of prompt injection defences often assumes that detector calibrations transfer across languages. This paper provides a rigorous empirical analysis of cross-lingual injection detection across 57,912 scored inferences spanning English, Chinese, Arabic, and Spanish. We show that multilingual safety classifiers do not suffer from cross-lingual recall collapse under English-calibrated thresholds. Rather, multilingual encoders exhibit pronounced calibration drift, assigning higher anomaly scores to non-English inputs regardless of intent. Under standard English operating points, benign false-positive rates escalate to 67.7\% for Chinese and 15.7\% for Arabic. Conversely, conditioning detection thresholds on language restores benign specificity but precipitates severe evasion collapse, reducing Arabic recall from 54.0\% to 7.0\%. Furthermore, agent execution evaluations over 1,160 episodes demonstrate an empirical floor effect where attack success remains near zero across languages, preventing valid claims of cross-lingual divergence in task execution. Finally, subword tokenisation fertility reaches 4.33x in non-Latin scripts, correlating with score inflation. We formalise these empirical trade-offs and provide methodological recommendations for evaluating multilingual security boundaries.
\end{abstract}

\section{Introduction}
Empirical NLP methods for safety benchmarking must establish robust measurement standards. In this paper, we examine how cross-lingual transfer impacts the empirical reliability of prompt injection detectors and agent execution safety.

\section{Empirical Evaluation Protocol}
We design an empirical testbed scoring:
\begin{enumerate}
    \item 1,160 InjecAgent execution episodes under open-weight models across four languages.
    \item 57,912 detector predictions across multilingual encoder families (e.g. DeBERTa v1, XLM-RoBERTa).
\end{enumerate}

\section{Empirical Observations}
Our findings show:
\begin{itemize}
    \item \textbf{Absence of Execution Divergence:} Floor effects dominate agent execution, with attack success rates bounded below 2\% across all languages.
    \item \textbf{Calibration Drift vs Evasion Collapse:} Under monolithic English thresholds, false alarm rates reach 67.7\% on Chinese benign text. Under language-specific calibration, Arabic detection recall drops by 47.0 percentage points.
\end{itemize}

\section{Tokenisation Fertility as an Empirical Factor}
We quantify subword fertility across typologically diverse scripts, demonstrating up to 4.33x expansion for non-Latin writing systems, and analyse its relationship to classifier score inflation.

\section{Limitations}
This empirical study evaluates open-weight architectures and existing benchmark corpora. Results reflect current public safety classifiers and may differ for proprietary systems with unknown multilingual training mixtures.

\section{Ethics Statement}
All evaluations rely on defensive research benchmarks and synthetic corpora. No active exploitation was conducted. The work supports improved calibration for multilingual NLP safety tools.

\bibliography{custom}

\end{document}
